% ch26b-cute-persist-d-flash-attn.tex — CuTe 应用（四）：sm_120 持久化 FlashAttention
% ——超越 cuDNN（2026-09-22，源码 flash_attn.cuh Phase 8，L3490-4170）
% 图：全 TikZ（用户指定）；代码段引用 ../flash_attn.cuh linerange。
\chapter{CuTe 应用（四）：sm\_120 持久化 FlashAttention——超越 cuDNN}\label{ch:26b}

\section{本章导读}

第~\ref{ch:25}~章结束时，我们手里已经有三个 CuTe 版 FlashAttention：cp.async 统一流水、
TMA + 单 consumer warp specialization、FA3 式双 consumer。它们回答的都是「布局事实搬进类型之后
attention 长什么样」，但性能口径上还欠一个交代：\textbf{这些教学实现离厂商库到底有多远？}
本章给出终局答案——把三个实现里各自最强的部分拼起来，再补上两个只有在「kernel 全权可控」时
才做得出的优化（softmax scale$\cdot\log_2 e$ 融合、persistent CTA），得到
\texttt{flash\_\allowbreak{}attn\_\allowbreak{}cute\_\allowbreak{}persist\_\allowbreak{}d\_\allowbreak{}sm120}
（flash\_attn.cuh L3620，\texttt{NOTES\_\allowbreak{}V2\_\allowbreak{}ENABLE\_\allowbreak{}CUTE}
宏块 Phase~8）：在 RTX PRO 5000（sm\_120，96 SM）上
$H{=}32$/$N{=}8192$/$D{=}128$ dense 场景实测 \textbf{236.9 TFLOPS，为 cuDNN SDPA
（227.5 TFLOPS）的 1.04 倍}。这是全书第一个（也是唯一一个）在绝对性能上越过厂商库的 kernel。

本章的四条主线，每条都对应「超越 cuDNN」的一个技术支柱：

\begin{itemize}[itemsep=1pt,topsep=2pt]
  \item \textbf{WS 双角色}（\S\ref{sec:26b-arch}）：128 线程 producer（只发 TMA）+
        256 线程 consumer（只做 MMA），\texttt{setmaxnreg} 把 producer 的寄存器
        让渡给 consumer——这是 sm\_120 上把「搬运」与「计算」物理解耦的唯一手段；
  \item \textbf{Q 寄存器持久化}（\S\ref{sec:26b-persistq}）：Q tile 对整个 KV 循环
        不变，smem$\to$寄存器只搬一次，之后每次 $QK^\top$ 都是 A-operand 常驻的
        \texttt{gemm\_rs}；
  \item \textbf{softmax scale 融合}（\S\ref{sec:26b-scale-math}）：把
        $s=\text{scale}\cdot\log_2 e$ 一次性乘进 Q fragment，softmax 的 max/exp
        两个循环里的逐元素乘法全部消失；
  \item \textbf{persistent CTA}（\S\ref{sec:26b-persistent}）：dense 场景 grid
        收缩到 SM 数，每个 CTA 串行吃多个 q-tile，消掉 wave 量化尾损。
\end{itemize}

\textbf{前置章节}：第~\ref{ch:15}--\ref{ch:16}~章（attention 数学与 FA2 online
softmax）、第~\ref{ch:17}~章（TMA 与 mbarrier 协议）、第~\ref{ch:21}--\ref{ch:22}~章
（TiledCopy 与 TiledMMA）、第~\ref{ch:25}~章（CuTe FA 三实现对照——本章直接复用其
五个工具函数）。\textbf{源码区间}：flash\_attn.cuh L3490--4170（Phase 8 宏块）。
\textbf{编译宏与架构}：\texttt{NOTES\_\allowbreak{}V2\_\allowbreak{}ENABLE\_\allowbreak{}CUTE}，
\textbf{编译目标必须 \texttt{sm\_120f}}（\texttt{sm\_120a} 会让 ptxas 以 C7506 静默丢弃
\texttt{setmaxnreg}，见 \S\ref{sec:26b-pitfalls}）。本章 kernel 只依赖 sm\_80 级的
mma.sync 与 sm\_90 级的 TMA/mbarrier，任何 Blackwell（sm\_120）显卡均可运行。

\section{问题与动机：为什么专用 kernel 能赢厂商库}\label{sec:26b-motivation}

「手写 kernel 赢 cuDNN」在 Hopper（sm\_90）时代几乎不可信：cuDNN 在 H100 上用
WGMMA + warp specialization + TMEM 调到了接近理论峰值，普通开发者用 CuTe 从零追赶
成本极高。但 sm\_120（Blackwell 消费线/专业线）改变了博弈格局——这块芯片\textbf{没有}
WGMMA（sm\_90 专属）、\textbf{没有} tcgen05/TMEM（sm\_100 专属），cuDNN 与我们一样，
只剩同一条张量核路径：

\begin{equation*}
  \text{mma.sync.aligned.m16n8k16.row.col.f32.f16.f16.f32}
\end{equation*}

也就是说，在 sm\_120 上「厂商库 vs 手写」的差距不再是指令代差，而全部落在\textbf{协议设计}
上：tile 几何、流水深度、寄存器分配、softmax 的指令数、wave 尾损。这五项恰恰是本书
前 25 章逐项打磨过的东西。实测基线（PRO 5000，$B{=}1$/$H{=}32$/$N{=}8192$/$D{=}128$
fp16，dense）：

\begin{center}
\begin{tabular}{lcc}
\toprule
实现 & TFLOPS & 相对 cuDNN \\
\midrule
cuDNN SDPA（\texttt{CUDNN\_ATTENTION} 后端） & 227.5 & 1.00$\times$ \\
本章 kernel & \textbf{236.9} & \textbf{1.04$\times$} \\
\bottomrule
\end{tabular}
\end{center}

领先从哪里来？逐项拆解（A/B 实测，独立 kernel 上各开关单独计量）：
\texttt{exp2f} 替代 \texttt{expf}（\texttt{-{}-use\_\allowbreak{}fast\_\allowbreak{}math}，
+2\%\textasciitilde4\%）、softmax scale 融合（+3\%\textasciitilde4\%）、persistent CTA
（dense 中等序列 +2\%\textasciitilde6\%）、FA-4 条件重缩放（长序列收益，复用
第~\ref{ch:16}~章结论）。没有任何一项是「黑魔法」，每一项都是把协议里的一次浪费
换成零开销。

还有一个工程动机：厂商库的 SDPA 对\textbf{非标准形状}（GQA 折叠、滑窗 causal、
head\_dim=96 这类非 2 的幂）经常落入慢速 fallback；专用 kernel 按自己的 tile 几何
原生支持这些形状——这是性能之外的第二条护城河。

\section{数学原理}\label{sec:26b-math}

\subsection{log2 域 online softmax 回顾}

第~\ref{ch:15}--\ref{ch:16}~章已经推导过 online softmax，这里只保留本章需要的形态。
对每个 KV tile，softmax 的在线更新为
\begin{align}
  m^{(t)} &= \max\!\big(m^{(t-1)},\ \max_j s_j^{(t)}\big), \label{eq:26b-max}\\
  \ell^{(t)} &= \ell^{(t-1)} e^{m^{(t-1)}-m^{(t)}} + \sum_j e^{s_j^{(t)}-m^{(t)}}, \label{eq:26b-sum}\\
  O^{(t)} &= O^{(t-1)} e^{m^{(t-1)}-m^{(t)}} + \sum_j e^{s_j^{(t)}-m^{(t)}} v_j^{(t)}, \label{eq:26b-o}
\end{align}
其中 $s_j = \text{scale}\cdot q\cdot k_j$。GPU 上把 $e^x$ 换成
$\mathrm{exp2}(x\cdot\log_2 e)$（一条 \texttt{MUFU.EX2} 即可），为此把
$\log_2 e$ 预乘进 scale：
\begin{equation}\label{eq:26b-log2domain}
  s_j^{\log_2} = \underbrace{\text{scale}\cdot\log_2 e}_{s}\cdot (q\cdot k_j),
  \qquad e^{s_j - m} = \mathrm{exp2}\!\big(s_j^{\log_2} - m^{\log_2}\big).
\end{equation}
此后所有 max/sum/rescale 都在 log2 域完成，epilogue 再乘回 $\ln 2$（若需要 lse）。

\subsection{FA-4 条件重缩放}

式~\eqref{eq:26b-sum}/\eqref{eq:26b-o} 里的 $e^{m^{(t-1)}-m^{(t)}}$ 要求每个 tile
都对 $\ell$ 和 $O$ 做一次逐元素乘法。FA-4 的观察是：随机数据下新 tile 的 max 很少
超过历史 max 太多。记 $\Delta = m^{(t-1)} - m^{(t)} \le 0$（log2 域），当
\begin{equation}\label{eq:26b-fa4}
  \Delta \ge -\tau, \qquad \tau = 8 = \log_2 256
\end{equation}
时，$2^{\Delta}$ 只会让 P 累计最多膨胀 $256\times$——fp16 的 P 最大 65504，
f32 累加器更无压力——因此可以\textbf{跳过}这次重缩放：$m$ 保持旧值
（stale max）、$O$/$\ell$ 原样累加。膨胀的部分在 epilogue 的 $1/\ell$ 里全部
抵消（分子分母同乘一个常数）。dense 数据下命中率高（实测 95\%+ 的 tile 满足
\eqref{eq:26b-fa4}），只有真的出现「更尖」的 tile 才付出重缩放代价。这正是第~\ref{ch:16}~章
手写版的结论，本章把它原样搬进 CuTe 版（\texttt{online\_\allowbreak{}safe\_\allowbreak{}softmax\_\allowbreak{}fa4}，
flash\_attn.cuh L3578）。

\subsection{softmax scale$\cdot\log_2 e$ 融合}\label{sec:26b-scale-math}

式~\eqref{eq:26b-log2domain} 里，$\log_2 e$ 折进 scale 只是「提前一次乘法」——
每个 score 元素仍然要乘 $s$ 两次：max pass 一次（\eqref{eq:26b-max}）、
exp pass 一次（\eqref{eq:26b-sum}）。对一个 $128\times 64$ 的 S tile，这是
$2\times 128\times 64 / 256\text{thread} = 64$ 次 FMUL/线程/tile，
而 QK GEMM 本身（m16n8k16 路径）每 tile 每线程只有 64 条 HMMA—— softmax 的
scale 乘法指令数与 tensor core 指令数同量级，是纯浪费。

本章的优化是把 $s$ 从「乘在 S 上」改为「乘在 Q 上」：
\begin{equation}\label{eq:26b-fold}
  S_{\text{fused}} = (s\odot Q)\,K^\top = s\cdot(QK^\top) = s\cdot S_{\text{raw}},
\end{equation}
其中 $\odot$ 是逐元素乘。等价性来自两条初等性质（$s>0$）：
\begin{align}
  \max_j (s\cdot x_j) &= s\cdot \max_j x_j, \label{eq:26b-maxcommute}\\
  \mathrm{exp2}(s\cdot x - m) &= \mathrm{exp2}(S_{\text{fused},j} - m). \label{eq:26b-expcommute}
\end{align}
Q fragment 在寄存器里、对整个 KV 循环不变（persist-D！），所以
\eqref{eq:26b-fold} 的代价只是 \texttt{size(tCrQ)} 次一次性 FMUL（Q s2r 后立即做），
外加 $s\odot Q$ 引入一次 fp16 重舍入。收益是 softmax 两个 pass 的逐元素乘法
\textbf{归零}。注意这条优化的前提是 persist-D：如果 Q 每个 tile 重新装载
（split-Q/FA3 布局），折叠本身也要每 tile 做一次，收益减半。

\begin{figure}[H]
\centering
\begin{tikzpicture}[
  font=\footnotesize,
  every node/.append style={align=center},
  box/.style={draw=tkBlue, fill=tkFillBlue, rounded corners=1.5pt, minimum height=0.7cm, inner sep=3pt},
  ops/.style={draw=tkPurple, fill=white, rounded corners=1.5pt, minimum height=0.7cm, inner sep=3pt},
  mul/.style={draw=tkOrange, fill=white, rounded corners=1.5pt, minimum height=0.6cm, inner sep=2pt, font=\scriptsize},
  arr/.style={-{Stealth[length=2.2mm]}, thick},
]
% before path
\begin{scope}
  \node[box] (q1) at (0,0) {$Q$};
  \node[box, right=0.5 of q1] (k1) {$K^\top$};
  \node[ops, right=0.5 of k1] (qk1) {$QK^\top$\\{\scriptsize HMMA}};
  \node[mul, right=0.35 of qk1] (m1a) {$\times s$};
  \node[mul, right=0.3 of m1a] (m1b) {$\times s$};
  \node[ops, right=0.35 of m1b] (sm1) {softmax};
  \draw[arr] (q1) -- (qk1);
  \draw[arr] (k1) -- (qk1);
  \draw[arr] (qk1) -- (m1a);
  \draw[arr] (m1a) -- node[above, font=\tiny] {max} (m1b);
  \draw[arr] (m1b) -- node[above, font=\tiny] {exp} (sm1);
  \node[font=\scriptsize\bfseries, above=0.15 of qk1] {before：每 tile 每 element 乘 2 次};
  \node[font=\tiny, tkOrange, below=0.12 of m1a] {$128{\times}64{\times}2$ FMUL / tile};
\end{scope}
% after path
\begin{scope}[shift={(0,-2.6)}]
  \node[mul] (fold) at (0,0) {$Q\times s$};
  \node[box, right=0.45 of fold] (k2) {$K^\top$};
  \node[ops, right=0.5 of k2] (qk2) {$(s\,Q)K^\top$\\{\scriptsize HMMA}};
  \node[ops, right=0.55 of qk2] (sm2) {softmax\\{\scriptsize 无乘法}};
  \draw[arr] (fold) -- (qk2);
  \draw[arr] (k2) -- (qk2);
  \draw[arr] (qk2) -- (sm2);
  \node[font=\scriptsize\bfseries, above=0.15 of qk2] {after：Q s2r 后一次性折叠};
  \node[font=\tiny, tkGreen, below=0.12 of fold] {$128{\times}D$ FMUL，全程仅一次};
\end{scope}
\end{tikzpicture}
\caption{softmax scale 融合的数据流对照。上：传统路径，S 从 QK GEMM 出来后在
softmax 的 max/exp 两个 pass 各乘一次 $s=\text{scale}\cdot\log_2 e$；
下：融合路径，$s$ 在 Q s2r 后一次性乘进 A-fragment（persist-D 使 Q 对整个
KV 循环不变），HMMA 直接输出 log2 域分数，softmax 循环体内零 FMUL。等价性由
式~\eqref{eq:26b-maxcommute}/\eqref{eq:26b-expcommute} 保证；代价是 $s\odot Q$
的一次 fp16 重舍入。}
\label{fig:26b-scale}
\end{figure}

\section{设计与实现}\label{sec:26b-design}

\subsection{总体架构：WS 双 warpgroup 与寄存器池数学}\label{sec:26b-arch}

kernel 用 384 线程、\texttt{\_\_launch\_\allowbreak{}bounds\_\allowbreak{}(384,~1)}：
前 128 线程（4 个 warp，1 个 warpgroup）是\textbf{ producer}，只做 TMA 发射与
mbarrier 维护；后 256 线程（8 个 warp，2 个 warpgroup）是\textbf{ consumer}，
只做 LDSM 装载、mma.sync、softmax。两个角色通过 K/V 的
full/empty mbarrier 对解耦（协议见 \S\ref{sec:26b-tma}）。

关键在寄存器：producer 只发 TMA（不碰 tensor core），128 个寄存器绰绰有余；
consumer 的 O 累加器（$128\times D$ f32/tile，每线程 $D/2{\times}4$ 个）+ P/S
fragment + softmax 状态却是寄存器大户。\texttt{setmaxnreg} 让两者在
SM 寄存器池（sm\_120 为 64K/SM）里重新分配：

\begin{equation}\label{eq:26b-regpool}
  \underbrace{4\,\text{warp}\times 32\,\text{T}\times 32}_{\text{producer }400}
  + \underbrace{8\,\text{warp}\times 32\,\text{T}\times 232}_{\text{consumer }2}
  = 4096 + 59392 = 63488 < 65536 .
\end{equation}

producer 在入口执行 \texttt{warpgroup\_\allowbreak{}reg\_\allowbreak{}dealloc<32>}
（让出 96 个/线程），consumer 执行 \texttt{warpgroup\_\allowbreak{}reg\_\allowbreak{}alloc<232>}
（拿满硬件上限的可用额度）。没有这一步，consumer 会被编译器压到人均 168 个寄存器，
O 累加器溢出到 local memory，性能直接腰斩。

\begin{figure}[H]
\centering
\begin{tikzpicture}[
  font=\footnotesize,
  wg/.style={draw, thick, rounded corners=2pt, minimum width=4.4cm, minimum height=1.5cm, inner sep=4pt, align=center},
  arr/.style={-{Stealth[length=2.2mm]}, thick},
]
% producer
\node[wg, draw=tkBlue, fill=tkFillBlue] (prod) at (0,0)
  {\textbf{producer}：128 T（4 warp）\\
   {\scriptsize TMA 发射（Q/K/V/O）+ mbarrier}\\
   {\scriptsize \texttt{reg\_dealloc<32>}：人均 32}};
% consumer
\node[wg, draw=tkGreen, fill=tkGreen!8] (cons) at (0,-3.2)
  {\textbf{consumer}：256 T（8 warp）\\
   {\scriptsize LDSM + mma.sync + softmax + rescale}\\
   {\scriptsize \texttt{reg\_alloc<232>}：人均 232}};
% smem
\node[wg, draw=tkPurple, fill=white, minimum width=2.6cm] (smem) at (5.6,-1.6)
  {SMEM\\{\scriptsize Q persist +}\\{\scriptsize K/V stage 池}};
% arrows
\draw[arr, tkBlue] (prod.east) -- node[above, font=\tiny, sloped] {TMA 写} (smem.north west);
\draw[arr, tkGreen] (smem.south west) -- node[above, font=\tiny, sloped] {LDSM 读} (cons.east);
\draw[arr, tkOrange, dashed] (cons.north) -- node[right, font=\tiny, align=left] {k\_empty/v\_empty\\（消费完成）} (prod.south);
\draw[arr, tkBlue, dashed] (prod.south west) ++(-1.2,0) -- node[left, font=\tiny, align=right] {k\_full/v\_full\\（数据就绪）} (cons.north west);
% register pool bar
\begin{scope}[shift={(-4.2,-1.6)}]
  \draw[very thick, tkRed] (0,1.6) -- (0,-1.6);
  \fill[tkBlue!25] (0,1.55) rectangle (0.5,1.15);
  \fill[tkGreen!30] (0,1.1) rectangle (0.9,-1.55);
  \node[font=\tiny, anchor=west] at (0.55,1.35) {4096};
  \node[font=\tiny, anchor=west] at (0.95,-0.2) {59392};
  \node[font=\tiny\bfseries, anchor=south, rotate=90] at (-0.25,0) {64K/SM};
\end{scope}
\end{tikzpicture}
\caption{WS 双角色架构与 \texttt{setmaxnreg} 寄存器再分配。左：SM 寄存器池按
式~\eqref{eq:26b-regpool} 重新切分，producer 只留 32/线程，consumer 拿 232/线程；
中：producer 只发 TMA、consumer 只算；右：smem 是两者的交接面。实线箭头是数据流
（TMA 写入 $\to$ LDSM 读出），虚线箭头是 mbarrier 控制流（full = 数据就绪、
empty = 消费完成），两者构成 \S\ref{sec:26b-tma} 的握手协议。}
\label{fig:26b-arch}
\end{figure}

\subsection{smem 布局与 swizzle}\label{sec:26b-smem}

smem 按 \texttt{[Q persist | K stages | V stages]} 一维排布，总量编译期由
\texttt{kSmemElems} 决定，D=128 时 $128{\times}128{\times}2 + 2{\times}(64{\times}128{\times}2)
\times 2 = 96$\,KB（opt-in 上限 99\,KB）。每段用哪个 swizzle atom 由\textbf{行字节数}
决定——这是第~\ref{ch:23}~章的结论在此处的落地：

\begin{lstinputlisting}[linerange={3530-3568},firstnumber=auto,
  title={flash\_attn.cuh L3530-3568（\texttt{Flash\allowbreak{}Attn\allowbreak{}Persist\allowbreak{}D\allowbreak{}Cu\allowbreak{}Te\allowbreak{}Traits}：smem 布局与双 TiledMma）}]{../flash_attn.cuh}
\end{lstinputlisting}

三个要点。第一，swizzle atom 按 $D\times 2$\,B 对 128/64 的整除性选 SW128/SW64，
保证 LDSM 读的每个 16\,B 段不跨 swizzle 周期。第二，\texttt{Smem\allowbreak{}Layout\allowbreak{}KVt}
不是新的存储，而是 \texttt{SmemLayoutKV} 与 col-major $(D, B_c)$ layout 的
\textbf{复合}（composition）——V 物理上仍按行存，PV GEMM 的 B-operand 却以转置
视图装载（\texttt{LDSM\_T}），免去 kernel 外物化 $V^\top$。第三，Q/K/V 共用
atom 但 tile 形状不同：Q 是 $128\times D$（持久），K/V 是 $B_c\times D$
（stage 轮转），$B_c$ 按 $D$ 反推 smem 预算（D=64 时 $B_c{=}128$、
D=128 时 $B_c{=}64$，注释里给了账）。

\begin{figure}[H]
\centering
\begin{tikzpicture}[font=\footnotesize, scale=0.96]
% smem bar
\fill[tkFillBlue]  (0,0)    rectangle (3.4,-1.5) node[midway, align=center] {\textbf{Q persist}\\{\scriptsize $128{\times}D$}\\{\scriptsize (epilogue 复用为 sO)}};
\fill[tkBlue!12]   (3.5,0)  rectangle (5.5,-1.5) node[midway, align=center] {\textbf{K[0]}\\{\scriptsize $B_c{\times}D$}};
\fill[tkBlue!12]   (5.6,0)  rectangle (7.6,-1.5) node[midway, align=center] {\textbf{K[1]}\\{\scriptsize $B_c{\times}D$}};
\fill[tkGreen!10]  (7.7,0)  rectangle (9.7,-1.5) node[midway, align=center] {\textbf{V[0]}\\{\scriptsize $B_c{\times}D$}};
\fill[tkGreen!10]  (9.8,0)  rectangle (11.8,-1.5) node[midway, align=center] {\textbf{V[1]}\\{\scriptsize $B_c{\times}D$}};
\draw[very thick] (0,0) rectangle (11.8,-1.5);
\draw (3.45,0) -- (3.45,-1.5);
\node[font=\scriptsize, anchor=west] at (0,-1.85) {D=128, $B_c{=}64$, $S_K{=}S_V{=}2$：$32 + 16{\times}4 = 96$\,KB $\le 99$\,KB};
% stage rotation
\begin{scope}[shift={(0.3,-3.6)}]
  \foreach \g/\s/\col in {0/0/tkBlue, 1/1/tkBlue, 2/0/tkBlue, 3/1/tkBlue, 4/0/tkBlue}{
    \node[draw=\col, thick, fill=white, minimum width=1.35cm, font=\tiny] at (\g*1.5,0) {$g{=}\g$};
    \node[font=\tiny, anchor=north] at (\g*1.5,-0.42) {stage $s{=}\s$};
    \node[font=\tiny, tkPurple, anchor=north] at (\g*1.5,-0.78) {phase $(\g/2)\&1$};
  }
  \draw[-{Stealth}, thick, tkOrange] (0.75,0.45) -- (6.15,0.45)
    node[midway, above, font=\tiny] {全局 kv 计数 $g$（跨 q-tile 连续累计）};
\end{scope}
\end{tikzpicture}
\caption{smem 一维布局（上）与 K/V stage 轮转（下）。Q 段在 KV 循环期间只读，
循环结束后复用为 O epilogue 的 STSM 暂存（\S\ref{sec:26b-epilogue}）；K/V 各自
独立的 $S{=}2$ stage 池按全局 kv 计数 $g$ 轮转：stage $=g\bmod S$、mbarrier
相位 $=(g/S)\&1$。相位用\textbf{全局}计数而不是每 q-tile 重置，是 persistent
CTA（\S\ref{sec:26b-persistent}）与普通版协议兼容的关键。}
\label{fig:26b-smem}
\end{figure}

\subsection{TMA 装载协议：descriptor、V-first 预取与 mbarrier 相位}\label{sec:26b-tma}

Q/K/V/O 四个方向的 gmem 访问全部走 TMA。BHND packed 的张量被压平成
$(B\cdot H\cdot N, D)$ 的 2D 坐标空间，batch/head 的行偏移用
\texttt{domain\_\allowbreak{}offset} 在 kernel 内动态叠加——这样四个方向只需要
各一个 TMA descriptor（\texttt{make\_\allowbreak{}tma\_\allowbreak{}copy} 在
launcher 里构建一次，\texttt{CUTLASS\_\allowbreak{}GRID\_\allowbreak{}CONSTANT}
放进 constant memory），无论 $B/H$ 多少都不用重建。GQA 的 head 映射
（$h \mapsto h/(H/H_{kv})$）在计算 \texttt{kv\_\allowbreak{}row\_\allowbreak{}offset}
时完成。

producer 的发射时序是本协议的精华：Q 先发（一次性全 $D$），然后
K/V 各预取 $S{-}1$ 个 tile 填满流水，主循环里每个 tile 领先 $S{-}1$ 发射下一个。
注意 \textbf{V-first-then-K} 的顺序——V 的消费点（PV GEMM）晚于 K 的消费点
（QK GEMM）整整一段 softmax，先发 V 让 K 的 TMA 传输填进 QK 计算的空隙：

\begin{figure}[H]
\centering
\begin{tikzpicture}[
  font=\scriptsize,
  lane/.style={minimum width=1.0cm, minimum height=0.55cm, inner sep=1pt, align=center},
  arr/.style={-{Stealth[length=1.8mm]}, thick},
]
\def\lw{7.6}
% time axis
\draw[-{Stealth}] (0,1.1) -- (\lw+2.2,1.1) node[right, font=\tiny] {时间};
\foreach \t/\lab in {0.6/$t_0$, 2.6/$t_1$, 4.6/$t_2$, 6.6/$t_3$}
  \draw (\t,1.0) -- (\t,1.2) node[above, font=\tiny] {\lab};
% producer lane
\node[anchor=east, font=\scriptsize\bfseries] at (-0.2,0.35) {producer};
\node[lane, draw=tkBlue, fill=tkFillBlue, minimum width=2.0cm] at (1.6,0.35) {Q TMA};
\node[lane, draw=tkBlue, fill=tkFillBlue] at (3.1,0.35) {K0};
\node[lane, draw=tkBlue, fill=tkFillBlue] at (4.15,0.35) {K1};
\node[lane, draw=tkPurple, fill=tkPurple!8] at (5.2,0.35) {V0};
\node[lane, draw=tkPurple, fill=tkPurple!8] at (6.25,0.35) {V1};
\node[lane, draw=tkBlue, fill=tkFillBlue] at (7.3,0.35) {K2};
\node[lane, draw=tkPurple, fill=tkPurple!8] at (8.35,0.35) {V2};
\node[lane, draw=tkBlue, fill=tkFillBlue] at (9.4,0.35) {K3};
% consumer lanes
\node[anchor=east, font=\scriptsize\bfseries] at (-0.2,-0.55) {consumer};
\node[lane, draw=tkGreen, fill=tkGreen!10] at (1.6,-0.55) {Q s2r $+\,s$ 折叠};
\node[lane, draw=tkOrange, fill=orange!10, minimum width=1.4cm] at (3.0,-0.55) {QK$_0$};
\node[lane, draw=gray, fill=gray!10, minimum width=0.8cm] at (4.0,-0.55) {SM$_0$};
\node[lane, draw=tkOrange, fill=orange!10, minimum width=1.4cm] at (5.0,-0.55) {QK$_1$};
\node[lane, draw=gray, fill=gray!10, minimum width=0.8cm] at (6.0,-0.55) {SM$_1$};
\node[lane, draw=orange!20, draw=tkOrange] at (7.2,-0.55) {PV$_0$};
\node[lane, draw=orange!20, draw=tkOrange] at (8.3,-0.55) {PV$_1$};
% dependency arrows
\draw[arr, tkBlue, dashed] (3.1,0.08) -- (3.0,-0.26) node[midway, right, font=\tiny] {k\_full[0]};
\draw[arr, tkPurple, dashed] (5.2,0.08) -- (7.2,-0.26);
\draw[arr, tkOrange, dashed] (3.75,-0.55) -- (3.75,-0.9) -- (4.35,-0.9) node[midway, below, font=\tiny] {k\_empty[0]：K1 的 TMA 可以复写 stage 0};
\end{tikzpicture}
\caption{producer/consumer 流水时序（$S_K{=}S_V{=}2$ 示意）。producer 每个 KV tile
领先 $S{-}1$ 发射下一个 tile 的 TMA，且 V 先于 K 发射（V 的消费点更晚，
K 的传输得以填进 QK 计算窗口）。虚线箭头是 mbarrier 握手：\texttt{k\_full[s]}
（事务屏障，TMA 完成 + expect\_tx 校验）放行 consumer 读取；
\texttt{k\_empty[s]}（consumer 每线程 arrive）放行 producer 复写该 stage。
S tile 的 softmax（SM$_i$）与下一 tile 的 QK 在 consumer 侧串行，但与
producer 侧的 TMA 完全重叠。}
\label{fig:26b-pipeline}
\end{figure}

mbarrier 相位是这套协议最容易写错的地方。本书第~\ref{ch:17}~章用的是
「每 tile 重置」的相对相位；persistent CTA 下 q-tile 循环不能每轮重置
（producer 可能领先 consumer 整整一个 q-tile），所以相位改由\textbf{全局}
kv 计数 $g$ 决定：

\begin{equation}\label{eq:26b-phase}
  \text{stage} = g \bmod S, \qquad \text{phase} = \lfloor g/S \rfloor \&\, 1 ,
\end{equation}

$g$ 在 producer 与 consumer 两侧各自按消费序列连续累计（跨 q-tile 不清零），
wait 与 arrive 的奇偶序自然配对。配套有一条铁律：\texttt{k\_\allowbreak{}empty}/\texttt{v\_\allowbreak{}empty}
的\textbf{初始} arrive 只在第一个 q-tile 开头做一次——之后每个 stage 的
「空」信号完全由上一轮消费 arrive 维持。若每轮 q-tile 都补一次初始 arrive，
mbarrier 的相位会被多推一格，producer 的 wait 永远等不到正确的奇偶
（compute-sanitizer 的 synccheck 会报 \emph{Barrier error: Missing wait}）。

\subsection{Q 寄存器持久化与 \texttt{gemm\_rs}：TiledMMA/TiledCopy 组合}\label{sec:26b-persistq}

FA2 的 split-Q 布局里 Q 每个 tile 都要从 smem 重装；persist-D 的观察是：
对固定 q-tile，$Q$ 在整个 KV 循环里不变，装载一次进寄存器后，
$S = QK^\top$ 的每次计算都是「A 在寄存器、B 在 smem」的形态——第~\ref{ch:22}~章
的 \texttt{gemm\_rs} 正是为此而生。装载与折叠发生在 q-tile 循环开头：

\begin{lstinputlisting}[linerange={3876-3890},firstnumber=auto,
  title={flash\_attn.cuh L3876-3890（Q s2r 一次 + scale$\cdot\log_2 e$ 折叠进 A-fragment）}]{../flash_attn.cuh}
\end{lstinputlisting}

\texttt{partition\_\allowbreak{}fragment\_\allowbreak{}A(sQ)} 按 QK 的
TiledMMA 切出每线程的 A-fragment；\texttt{retile\_D} 把 copy 的目的视图对齐到
同一片寄存器（第~\ref{ch:21}~章的 TV-layout 分区）。折叠循环
\texttt{tCrQ(i) = (Element)(float(tCrQ(i)) * (scale * FA4\_LOG2E))} 逐元素乘
$s$ 后写回 fp16——此后 QK 的 HMMA 输出\textbf{直接}就是 log2 域分数
（式~\eqref{eq:26b-fold}），\S\ref{sec:26b-scale-math} 的全部收益在此兑现。

softmax 本体是 \texttt{online\_\allowbreak{}safe\_\allowbreak{}softmax\_\allowbreak{}fa4}
（L3578--3617）：与第~\ref{ch:16}~章手写版同构，差别只有两处——行/列视图来自
\texttt{convert\_\allowbreak{}layout\_\allowbreak{}acc\_\allowbreak{}rowcol}
（m16n8k16 C-fragment 的行列重解释，第~\ref{ch:22}~章），以及 scores 不再乘
scale（已经折叠）。行内 4-lane 蝴蝶归约（\texttt{\_\_shfl\_xor} 1 与 2）把
散在 4 个 lane 的同 row 分数归到一起——这是 m16n8k16 fragment 约定的直接
推论（每行 8 个分数由 4 个 lane 各持 2 个）。

PV 的 A-operand 复用第~\ref{ch:25}~章讲透的寄存器重解释技巧：
softmax 后的 P（f32 累加器 fragment）经 \texttt{convert\_\allowbreak{}type}
压成 fp16，再经 \texttt{convert\_\allowbreak{}layout\_\allowbreak{}acc\_\allowbreak{}Aregs<TiledMmaPV>}
原地重排成 PV 的 A-operand 布局——P 不落 smem，寄存器直通下一个 MMA。
V 侧则用 \S\ref{sec:26b-smem} 的转置视图 + \texttt{LDSM\_T} 装载。

\subsection{persistent CTA：wave 量化尾损与线性分配}\label{sec:26b-persistent}

非 persistent launch 的 grid $=$ q-tile 总数，硬件按 wave 调度：
96 个 SM 一波。若 tile 数是 $96k + r$（$0<r<96$），最后一个 wave 只有
$r$ 个 SM 在干活，其余空转——这就是 wave 量化尾损。$H{=}32$、
$N{=}8192$、$k_B{=}128$ 时 tile 数 $= 32\times 64 = 2048 = 21\times 96 + 32$，
最后一个 wave 尾损约 $\tfrac{96-32}{96\times 22} \approx 3\%$。

persistent CTA 的做法：grid 收缩到 $\min(\text{tiles},\ \text{SM 数})$，
每个 CTA 用 \texttt{for (cta = blockIdx.x; cta < total; cta += gridDim.x)}
线性认领多个 q-tile。尾损变成「每个 CTA 多干一小段」，均匀分摊。
协议上有两个配套改动：其一就是 \S\ref{sec:26b-tma} 的全局 kv 计数相位
（式~\eqref{eq:26b-phase}）；其二是 \texttt{epi\_\allowbreak{}done} 屏障——
Q smem 在 epilogue 里被复用为 O 的 STSM 暂存，producer 想装下一个 q-tile
的 Q 必须等 consumer 把 O 写完，\texttt{epi\_\allowbreak{}done}（256 线程
arrive 的 CTA 屏障）就是这张「Q smem 空了」的回执。

\begin{figure}[H]
\centering
\begin{tikzpicture}[font=\scriptsize, scale=0.9]
% non-persistent
\begin{scope}
  \foreach \w in {0,...,4}{
    \foreach \s in {0,...,9}{
      \pgfmathtruncatemacro{\last}{\w==4 ? 3 : 9}
      \ifnum\s>\last\fill[gray!8] (\s*0.62,\w*0.5) rectangle (\s*0.62+0.62,\w*0.5+0.5);
      \else\fill[tkBlue!30] (\s*0.62,\w*0.5) rectangle (\s*0.62+0.62,\w*0.5+0.5);
      \fi
    }
  }
  \draw[very thick] (0,0) rectangle (6.2,2.5);
  \draw[tkRed, very thick, dashed] (2.48,2.0) rectangle (6.2,2.5);
  \node[font=\tiny, anchor=north] at (3.1,2.55) {wave 1（10 tile/CTA）};
  \node[font=\tiny, anchor=north] at (3.1,-0.08) {wave 3};
  \node[font=\tiny, tkRed, align=center, anchor=west] at (6.4,2.25)
    {wave 4：只有 4 个\\SM 有活\\{\bfseries 尾损 $\tfrac{6}{44}\approx 14\%$}};
\end{scope}
% persistent
\begin{scope}[shift={(0,-4.6)}]
  \foreach \w in {0,...,3}{
    \foreach \s in {0,...,9}{
      \pgfmathtruncatemacro{\tail}{\w==3 ? 4 : 9}
      \ifnum\s>\tail\fill[gray!8] (\s*0.62,\w*0.5) rectangle (\s*0.62+0.62,\w*0.5+0.5);
      \else\fill[tkGreen!25] (\s*0.62,\w*0.5) rectangle (\s*0.62+0.62,\w*0.5+0.5);
      \fi
    }
  }
  \draw[very thick] (0,0) rectangle (6.2,2.0);
  \node[font=\tiny, anchor=north] at (3.1,2.05) {每 CTA 4.4 个 tile：尾损 $\tfrac{0.4}{17.6}\approx 2\%$};
\end{scope}
\node[font=\footnotesize\bfseries, anchor=south] at (3.1,2.85) {非 persistent：grid = 44 tile};
\node[font=\footnotesize\bfseries, anchor=south] at (3.1,-2.3) {persistent：grid = 10 SM};
\end{tikzpicture}
\caption{wave 量化尾损示意（10 个 SM、44 个 tile 的缩小模型）。上：非
persistent launch 按整 wave 调度，最后一个 wave 只有 4 个 SM 在算，其余
空转满一个 tile 时长；下：persistent CTA 把 grid 收缩到 SM 数，每个 CTA
串行吃 $\lceil 44/10\rceil$ 或 $\lfloor 44/10\rfloor$ 个 tile，不均衡被压缩到
最后一个 tile 的零头。真实场景（96 SM，$H{=}32$/$N{=}8192$）dense 收益
+2\%\textasciitilde6\%；\textbf{causal 反向}——tile 负载前轻后重，线性分配会把
重 tile 锁死在同一批 CTA 上（实测 $-9\%\sim-21\%$），此时必须回退全量 grid
让硬件调度器自然均衡。}
\label{fig:26b-persistent}
\end{figure}

dense 与 causal 的分派就在 launcher 里一行：
\begin{lstinputlisting}[linerange={4130-4148},firstnumber=auto,
  title={flash\_attn.cuh L4130-4148（launcher 的 grid 策略与 smem opt-in）}]{../flash_attn.cuh}
\end{lstinputlisting}

\subsection{Epilogue：R$\to$S(STSM)$\to$TMA store 与尾部 R$\to$G}\label{sec:26b-epilogue}

KV 循环结束后，O 累加器除以 $\ell$（FA-4 的膨胀在此全部抵消），然后按 tile
形状分两条写出路径。整 tile：寄存器 $\to$ smem（\texttt{STSM}，复用 Q 段——
它已经完成历史使命）$\to$ TMA store；尾部 tile（$N_q \bmod 128 \ne 0$）：
逐元素行 guard 直接 R$\to$G 写 gmem。这里有个致命细节：producer 已经
\texttt{return} 退出，任何 \texttt{\_\_syncthreads} 都会死锁——CTA 级同步必须
换成 \texttt{Named\allowbreak{}Barrier::\allowbreak{}sync(256,~0)}（仅 consumer 参与）。
写完 O 才 \texttt{arrive} \texttt{epi\_\allowbreak{}done}，producer 由此获准复写
Q smem 装下一个 q-tile。

\begin{figure}[H]
\centering
\begin{tikzpicture}[
  font=\footnotesize,
  box/.style={draw, thick, rounded corners=1.5pt, minimum height=0.8cm, inner sep=3pt, align=center},
  arr/.style={-{Stealth[length=2.2mm]}, thick},
]
% full tile path
\begin{scope}
  \node[box, draw=tkGreen, fill=tkGreen!10] (reg) at (0,0) {O 寄存器\\{\scriptsize $\div\,\ell$、fp16}};
  \node[box, draw=tkBlue, fill=tkFillBlue] (smem) at (2.9,0) {sO（Q 段复用）\\{\scriptsize STSM}};
  \node[box, draw=tkPurple, fill=white] (gmem) at (5.9,0) {gmem O\\{\scriptsize TMA store}};
  \draw[arr] (reg) -- node[above, font=\tiny] {R$\to$S} (smem);
  \draw[arr] (smem) -- node[above, font=\tiny] {S$\to$G} (gmem);
  \node[font=\scriptsize\bfseries, above=0.1 of reg] {整 tile（$N_q \bmod 128 = 0$）};
\end{scope}
% tail path
\begin{scope}[shift={(0,-2.2)}]
  \node[box, draw=tkGreen, fill=tkGreen!10] (reg2) at (0,0) {O 寄存器\\{\scriptsize $\div\,\ell$、fp16}};
  \node[box, draw=tkOrange, fill=orange!8] (g2) at (3.6,0) {gmem O\\{\scriptsize 逐元素 row guard}};
  \draw[arr] (reg2) -- node[above, font=\tiny] {R$\to$G（行 $<N_q$ 才写）} (g2);
  \node[font=\scriptsize\bfseries, above=0.1 of reg2] {尾部 tile};
\end{scope}
% barrier note
\node[draw=tkRed, thick, rounded corners=2pt, font=\scriptsize, align=center, anchor=west] at (7.6,-1.1)
  {producer 已退出\\$\Rightarrow$ \texttt{\_\_syncthreads} 死锁\\$\Rightarrow$ \texttt{NamedBarrier(256)}\\写完 \texttt{arrive(epi\_done)}};
\end{tikzpicture}
\caption{Epilogue 双路径。整 tile 走 R$\to$S$\to$G：STSM 把 O fragment 写进
（已完成使命的）Q smem 段，NamedBarrier 后由一个线程发射 TMA store——
与装载对称的「反向 TMA 全链」。尾部 tile 行数不足 128，STSM/TMA 的整块写
会越界，退化为带行 guard 的 R$\to$G。}
\label{fig:26b-epilogue}
\end{figure}

\section{性能分析}\label{sec:26b-perf}

主场景实测（RTX PRO 5000，96 SM，fp16，min-of-30 events）：$B{=}1$/$H{=}32$/
$N_{kv}{=}8192$/$D{=}128$。TFLOPS 按 $4BH\,N_qN_{kv}D$（causal 减半）计。

\begin{center}
\begin{tabular}{lcccc}
\toprule
场景 & 本章 kernel & cuDNN SDPA & 比值 & 备注 \\
\midrule
dense self & \textbf{236.9} & 227.5 & \textbf{1.04$\times$} & 主对照 \\
causal & 208.9 & 202.1 & 1.03$\times$ & persistent 关闭（全量 grid） \\
GQA ($H_{kv}{=}8$) & 236.1 & 218.3 & 1.08$\times$ & cuDNN 小 head 数吃亏 \\
cross ($N_q{=}1024$) & 180.7 & 165.5 & 1.09$\times$ & 短 $N_q$ 也领先 \\
$N{=}16384$ dense & 243.2 & 236.9 & 1.03$\times$ & 大 $N$ 稳定 \\
\bottomrule
\end{tabular}
\end{center}

优化贡献的逐项拆解（同一 kernel 上各开关单独 A/B）：fast\_math
（\texttt{exp2f}）+2\%\textasciitilde4\%；scale 融合 +3\%\textasciitilde4\%（\S\ref{sec:26b-scale-math}，
指令数视角：每 tile 每线程省 64 条 FMUL，与 HMMA 数同量级）；
persistent CTA dense 中等序列 +2\%\textasciitilde6\%、causal $-9\%\sim-21\%$（分支
关闭）；FA-4 条件重缩放 dense 命中率 95\%+，与第~\ref{ch:16}~章结论一致。
各项近似可加，合计与 1.04$\times$ 的总领先吻合。

与第~\ref{ch:25}~章三个实现的纵向对比更能说明「协议设计」的价值：
实现 1（cp.async 统一流水）受限于装载与计算串行；实现 2（TMA+WS 单 consumer）
解耦了搬运但 Q 每 tile 重装、softmax 带逐元素乘法、wave 尾损裸奔；
本章在实现 2 的骨架上叠满四项优化，同一块 GPU 上从「勉强追平 SDPA」
推进到「稳定 +4\%」。

\section{勘误与考据}\label{sec:26b-kaogu}

本章源码为 2026-09-22 新增（Phase 8），注释与正文同批撰写，无历史勘误。
生产对照版（ffpa-attn 的 \texttt{persist\_\allowbreak{}d\_\allowbreak{}ws\_\allowbreak{}fwd\_\allowbreak{}cute\_\allowbreak{}sm120}）
与本章 kernel 同源，其 RFC（PC-18/PC-19）记录了 scale 融合与 persistent CTA
的生产 A/B 数据，本章表格与之交叉一致。

\section{常见坑}\label{sec:26b-pitfalls}

\begin{enumerate}[itemsep=2pt]
  \item \textbf{\texttt{sm\_120a} 静默丢弃 \texttt{setmaxnreg}}：ptxas 对 120a
        报 C7506 后直接把指令扔掉——kernel 结果正确、性能腰斩，且\textbf{无任何
        运行时告警}。必须编译 \texttt{sm\_120f}。诊断手段：\texttt{cuobjdump
        -{}-dump-sass} 搜 \texttt{SETMAXNREG}。
  \item \textbf{persistent 初始 arrive 必须只做一次}：每 q-tile 都补
        \texttt{k\_empty} 初始 arrive 会把 mbarrier 相位多推一格，producer 的
        wait 奇偶错位（compute-sanitizer synccheck 报 \emph{Missing wait}）。
        相位改由全局计数 $g$ 决定（式~\eqref{eq:26b-phase}）后，「空」信号由
        消费 arrive 自持。
  \item \textbf{producer 退出后禁用 \texttt{\_\_syncthreads}}：CTA 内任何
        barrier 都必须换成仅 consumer 参与的
        \texttt{NamedBarrier}（256 线程），否则 epilogue 死锁。
  \item \textbf{gridDim 携带大数撞 65535 上限}：早期版本把
        $B\cdot H\cdot N$ 塞进 \texttt{gridDim.z}，$N_q{=}4096$、$H{=}16$ 时
        恰好 $65536 > 65535$ 报 \texttt{cudaErrorInvalidValue}——行数改走
        kernel 参数。
  \item \textbf{依赖类型少 \texttt{typename}}：kernel 模板参数
        \texttt{Traits} 是依赖类型，直接写
        \texttt{Traits::SmemCopyAtom\{\}} 编不过；先
        \texttt{using SmemCopyAtom =} \texttt{typename Traits::SmemCopyAtom;}
        再用。第~\ref{ch:25}~章（具体类型别名）没有这个问题，本章踩过。
\end{enumerate}

\section{面试要点速查}\label{sec:26b-interview}

\begin{itemize}[itemsep=1pt]
  \item \textbf{Q：为什么你的 kernel 能赢 cuDNN？}\\
        A：sm\_120 没有 WGMMA/TMEM，大家都在 mma.sync 上；差距全在协议——
        Q 持久化省装载、scale 融合省 softmax 指令、persistent 省 wave 尾损、
        WS+TMA 流水藏延迟。四项合计 +4\%。
  \item \textbf{Q：scale 融合的数学依据与代价？}\\
        A：$\max(sx)=s\max(x)$（$s>0$）与 exp2 线性平移；$S_{\text{fused}}=(sQ)K^\top$。
        代价是 $s\odot Q$ 一次 fp16 重舍入（尾数 11 bit）；收益是每 tile 每
        线程 64 条 FMUL 归零。前提是 persist-D（Q 不变）。
  \item \textbf{Q：persistent CTA 为什么 causal 反而慢？}\\
        A：causal 的 tile 负载前轻后重，线性认领会把重 tail 锁死在同一批
        CTA；全量 grid 让硬件调度器填空。负载均匀才适合 persistent。
  \item \textbf{Q：\texttt{setmaxnreg} 的池数学？}\\
        A：$4w\times 32T\times 32 + 8w\times 32T\times 232 = 63488 < 64K$；
        producer 让出的正是 consumer 多拿的。无此指令则编译器均分 168/线程，
        O 累加器溢出。
  \item \textbf{Q：mbarrier 相位为什么用全局计数？}\\
        A：persistent 下 producer 可能跨 q-tile 领先消费，每 tile 重置相位
        无法与全局发射序配对；$g$ 连续累计 + $\lfloor g/S\rfloor\&1$ 保持
        wait/arrive 奇偶一致。
\end{itemize}

\section{最小测试与动手实验}\label{sec:26b-test}

本章测试在 \texttt{book/tests/ch26b\_\allowbreak{}fa\_\allowbreak{}persist\_\allowbreak{}d.cu}：
CPU fp64 参考对照，容差 \texttt{TOL\_F16ACC}（$5\times 10^{-2}$），十个 case 覆盖
dense / causal / GQA / Q 尾部 tile（$N_q{=}300$ 的 R$\to$G 路径）/
persistent 多迭代（$H{=}32$、$N{=}512$ 时 128 个 tile $>$ 96 SM，触发
\texttt{epi\_done} 与全局 kv 计数路径）/ KV 尾 mask（$N_{kv}{=}300$ 非
$k_{Bc}$ 倍数）/ causal 正反滑窗（$N_{kv}{>}N_q$ 与 $N_{kv}{<}N_q$，后者
前段 q 行全 mask，验证 $T_c^{eff}$ 钳零与 $L{=}0$ 行输出 0）/ $D{=}96$：

\begin{lstlisting}[style=console]
nvcc -std=c++20 -O2 -gencode arch=compute_120f,code=sm_120f \
  -DNOTES_V2_ENABLE_CUTE -I ../../third-party/cutlass/include \
  book/tests/ch26b_fa_persist_d.cu -o /tmp/ch26b_test && /tmp/ch26b_test
# PASS fa persist-D dense   D=64  B=1 H=2/2   Nq=256 Nkv=256 ...
# PASS fa persist-D dense   D=128 B=1 H=2/2   Nq=256 Nkv=256 ...
# PASS fa persist-D causal  D=128 B=1 H=2/2   Nq=512 Nkv=512 ...
# PASS fa persist-D dense   D=64  B=1 H=4/2   Nq=256 Nkv=256 ...
# PASS fa persist-D dense   D=128 B=1 H=2/2   Nq=300 Nkv=512 ...
# PASS fa persist-D dense   D=64  B=1 H=32/32 Nq=512 Nkv=512 ...
# PASS fa persist-D dense   D=128 B=1 H=2/2   Nq=256 Nkv=300 ...
# PASS fa persist-D causal  D=64  B=1 H=2/2   Nq=256 Nkv=512 ...
# PASS fa persist-D causal  D=128 B=1 H=2/2   Nq=512 Nkv=256 ...
# PASS fa persist-D dense   D=96  B=1 H=2/2   Nq=256 Nkv=256 ...
# ALL OK
\end{lstlisting}

集成入口在 notes-v2.cu（\texttt{test\_\allowbreak{}flash\_\allowbreak{}attn\_\allowbreak{}cute\_\allowbreak{}persist\_\allowbreak{}d\_\allowbreak{}sm120}，
含 $N_q{=}2048$/$H{=}8$ 的多 q-tile 大 case）；性能复测脚本
\texttt{bench\_pd.cu}（位于 \texttt{book/.tmp/ch26b/}；dense、causal、GQA、
$N$ 扫描，数据落 \texttt{.tmp/book-bench/ch26b/}）。动手实验建议：
(1) 把 \texttt{FA4\_\allowbreak{}RESCALE\_\allowbreak{}THRESHOLD} 调成 0
观察 FA-4 跳过率归零后的性能退化；(2) 强制 dense 也用全量 grid，
验证 \S\ref{sec:26b-persistent} 的 wave 尾损数字；(3) 去掉 Q s2r 后的
折叠循环，把 scale 乘法放回 softmax，复测 scale 融合的贡献。

\section*{延伸阅读}

\begin{itemize}[itemsep=1pt]
  \item FlashAttention-3 论文 \S3.1.4（conditional rescaling 的原始定义与
        $\tau$ 的选取依据）；
  \item CUTLASS 77\_blackwell\_fmha（FA-4 式 softmax 在 Blackwell 上的
        warp-uniform rescale 与 TMEM 形态——对照本章 per-row 形态）；
  \item @DefTruth《图解：从 Online-Softmax 到 FlashAttention V1/V2/V3》
        （online softmax 推导基线，本书第~\ref{ch:15}--\ref{ch:16}~章同源）；
  \item ffpa-attn RFC PC-18/PC-19（scale 融合与 persistent CTA 的生产版
        A/B 数据与量化家族适用性盘点）。
\end{itemize}
